%% kapa: exam=1 round=23 year=2012 period=1 form=A subject=civil_law subject_ko=민법 seq=8 number=8
%% source: Q-Net 감정평가사 기출문제 (cst00309) artlSeq=5052578 file=raw/1008/5052578/2012년 제23회 감정평가사 제1차시험 기출문제(1교시 A형).pdf
%% conversion: kordoc 4.24.51 -> tools/gen1.py (auto); review=proofread-sonnet; flags=none
\begin{question}{8}
전세권에 관한 설명으로 옳은 것은? (다툼이 있으면 판례에 의함)
\begin{choices}
  \item 전전세권의 설정에는 원전세권자와 전전세권자 사이의 전전세권 설정의 합의 및 등기와 원전세권설정자의 동의를 요한다.
  \item 구분등기 되지 않은 건물 일부의 전세권자는 전세권의 목적물이 아닌 나머지 건물부분에 대해서는 전세권에 기하여 경매신청을 할 수 없으나 우선변제권은 가진다.
  \item 전세권이 법정 갱신된 경우 갱신에 따른 등기를 하지 않으면, 전세권자는 그 목적물을 취득한 제3자에 대하여 갱신된 권리를 주장하지 못한다.
  \item 전세금을 현실적으로 지급하지 아니하고 기존의 채권으로 갈음하기로 한 경우에는 전세권이 성립 될 수 없다.
  \item 전세권 존속기간이 만료된 경우 전세권의 용익물권적 권능은 소멸되고 전세금반환채권은 무담보채권으로 전환된다.
\end{choices}
\end{question}
\answer{2}
